\documentclass[12pt]{article}
\usepackage[T1]{fontenc}
\usepackage[utf8]{inputenc}
\usepackage{lmodern}
\usepackage{textcomp}
\usepackage{enumitem}
\usepackage{geometry}
\usepackage{graphicx}
\usepackage{amsmath, amssymb}
\geometry{margin=1in}
\begin{document}
\begin{center}
Economics - Graduate Level Assessment 3 \
Total Marks: 40 \
Answer all questions.
\end{center}

\section*{Multiple Choice (10 marks)}
\begin{enumerate}[label=\arabic*.]
\item Suppose that there are only two goods: x and y. Which of the following is NOT correct?
\begin{enumerate}[label=(\alph*)]
    \item One can have comparative advantage in producing both goods.
    \item One can have both an absolute advantage and a comparative advantage in producing x.
    \item One can have absolute advantage and no comparative advantage in producing x.
    \item One can have comparative advantage and no absolute advantage in producing x.
\end{enumerate}
\item The primary focus of microeconomics is
\begin{enumerate}[label=(\alph*)]
    \item the overall national economy and its growth.
    \item the impact of government regulations on businesses.
    \item the stock market and its fluctuations.
    \item the economic impact of natural disasters.
    \item government.
\end{enumerate}
\item Which of the following is true in the long run in perfect competition?
\begin{enumerate}[label=(\alph*)]
    \item P = MR = MC = ATC
    \item P = MR = MC \textgreater{} ATC
    \item P \textgreater{} MR = MC = ATC
    \item P = MR \textgreater{} MC = ATC
\end{enumerate}
\item If technology makes production less expensive and at the same time exports decrease which of the following will result with certainty?
\begin{enumerate}[label=(\alph*)]
    \item Real GDP will decrease.
    \item The unemployment rate will increase.
    \item The country's trade balance will worsen.
    \item The interest rate will decrease.
    \item The price level will increase.
\end{enumerate}
\item Consider a series that follows an MA(1) with zero mean and a moving average coefficient of 0.4. What is the value of the autocovariance at lag 1?
\begin{enumerate}[label=(\alph*)]
    \item 1
    \item 2
    \item -0.4
    \item The autocovariance at lag 1 is undefined for an MA(1) series.
    \item 0.2
\end{enumerate}
\end{enumerate}

\section*{True / False (10 marks)}
\textit{Answer True or False}
\begin{enumerate}[label=\arabic*.]
\setcounter{enumi}{5}
\item True or False: Increasing government spending without changing taxes can stimulate real GDP growth while minimizing inflationary pressure.
\item True or False: Economic stabilization policy operates without any time-lags, allowing immediate effects on economic conditions.
\item True or False: An autonomous increase in spending refers to total spending that occurs independently of any changes in income levels.
\item True or False: A decrease in the number of firms in a market will lead to a rightward shift in the labor demand curve.
\item True or False: Higher unemployment rates in the United States typically lead to the appreciation of the U.S. dollar relative to the Chinese yuan.
\end{enumerate}

\section*{Long Form Response (20 marks)}
\textit{Answer each question with a clear and complete explanation.}
\begin{enumerate}[label=\arabic*.]
\setcounter{enumi}{10}
\item Discuss the significance of the principle of diminishing returns in growth theory, particularly in relation to capital accumulation and population growth's impact on output.
\item Discuss how to calculate the reserve ratio for a bank with $500 in checking deposits and a required reserve of $50. Additionally, analyze the implications of having \$475 in excess reserves.
\end{enumerate}

\vfill
\noindent\textit{End of Paper}
\end{document}